\chapter{Discussion et conditions d’interprétation}\label{ch:discussion}

\section{Un corpus daté ne résout pas à lui seul la chronologie}

Les inscriptions datées sont utiles pour comparer des formes et des usages, à condition que l’événement daté et le texte observé appartiennent à une même phase documentable. Un document qui associe une naissance et une déposition n’offre pas deux années également pertinentes pour dater son formulaire ; une plaque préparée du vivant peut avoir reçu une addition ; une dalle qui commémore plusieurs défunts porte plusieurs dates (chapitre~\ref{ch:evenements}). Le retour aux éditions et au support conditionne l’emploi des dates comme jalons stylistiques.

La table des fastes utilisée pour identifier les consuls n’a pas été collationnée sur les \emph{Consuls of the Later Roman Empire} : elle suit la tradition des fastes telle qu’elle est reçue dans les bases, et l’EDB, qui cite cet instrument, a servi de contrôle indirect. Les divergences d’itération, de collègue et de post-consulat restent ouvertes lorsqu’aucun témoin contrôlé ne permet de choisir. La vraisemblance d’une correction ne vaut pas preuve de la leçon.

\section{Deux bases ne font pas une édition}

La confrontation de l’EDCS et de l’EDB a été le principal instrument de cette révision. Elle a permis de récupérer des formules de date, des lignes omises, des parties grecques et des données matérielles ; elle a révélé des erreurs de saisie et des restitutions divergentes. Mais deux bases qui concordent ne garantissent pas une lecture : elles suivent souvent la même édition, et leurs accords peuvent reproduire la même erreur. Les « lectures stabilisées » de ce volume sont des lectures concordantes entre deux témoins secondaires ; elles ne sont pas des lectures collationnées sur les ICVR ni sur les monuments. Un export est un état de transmission, non un substitut transparent à l’édition.

\section{Le formulaire n’est pas une confession de foi exhaustive}

La présence d’une formule de paix ou de repos renseigne sur un choix de commémoration ; elle ne fournit pas une échelle d’adhésion doctrinale. L’absence d’une énumération du Père, du Fils et de l’Esprit dans \emph{ancilla dei et Chr(isti)} ne contredit pas une croyance trinitaire ; la présence de \emph{trinitatem} au \nnum{474} n’autorise pas à reconstruire des lignes perdues qui nommeraient les trois personnes.

L’énoncé fragmentaire et le genre du texte sont deux contraintes distinctes. Un texte bref a moins d’espace pour développer une doctrine ; un texte lacunaire peut avoir perdu ce qui la formulait. Sans série comparative définie et sans contrôle des longueurs et des pertes, le corpus ne permet pas de répartir les causes de la rareté du langage doctrinal (chapitre~\ref{ch:rarete}).

Le signe figuré pose un problème de même nature. Le chrisme, présent sur un texte sur sept parmi ceux que l’EDB décrit, dit l’appartenance chrétienne du monument plus souvent que les mots. Sa description dépend toutefois des rédacteurs de la base, et sa fonction — compendium du nom, signe juxtaposé, ornement — ne peut être établie que sur la pierre ou sur une photographie.

\section{Familles, opérateurs et institutions}

Les achats d’emplacements auprès de fossoyeurs attestent des opérations et des acteurs nommés. Ils ne suffisent pas à établir la propriété juridique des cimetières ni le degré de centralisation de leur gestion. La mention d’un évêque n’en fait pas automatiquement le vendeur : le datif \emph{Pascasio ep(is)c(opo)} du \nnum{400} appelle une analyse de relation et de fonction. La restitution \emph{a presbyteris} du \nnum{387} ne peut être mise sur le même plan qu’un nom de fonction lu. De même, la proximité recherchée des saints, telle que le \nnum{232} la donne à lire, ne mesure pas la distribution sociale des emplacements. Ces dossiers peuvent être confrontés aux problématiques de Charles Pietri, d’Éric Rebillard et de Steffen Diefenbach ; les développements de ces ouvrages qui ne sont connus ici que par leurs notices ou par des comptes rendus ne sont pas invoqués comme arbitres d’une décision.

\section{Les données numériques comme témoins}

La lecture méthodologique de Heřmánková, Kaše et Sobotková invite à distinguer le classement religieux de la fonction sociale, et la précision apparente d’une donnée de sa résolution réelle\footnote{P.~Heřmánková, V.~Kaše, A.~Sobotková, « Inscriptions as Data: Digital Epigraphy in Macro-historical Perspective », \emph{Journal of Digital History}, 1/1, 2021, DOI 10.1515/jdh-2021-1004.}. Dans ce volume, cette exigence se traduit par trois règles : ne pas confondre fiches et monuments ; ne pas confondre coordonnées bibliographiques et contexte archéologique ; ne pas confondre proportions documentaires et fréquences de population.

\section{Ce qui reste à faire}

Trois chantiers conditionnent l’exploitation historique de ce corpus au-delà de ce que le présent volume établit.

\begin{enumerate}
\item \emph{La collation des éditions.} Les textes à divergence forte (\st{edbDiv}), les seize textes absents de l’EDB et les formules de date contestées doivent être collationnés sur les ICVR et, pour les textes antérieurs, sur les éditions de De Rossi et de Diehl.
\item \emph{Le contrôle photographique.} Les reproductions signalées par l’EDB (\st{nFacsim} textes) permettraient de trancher la question des monogrammes, de vérifier les lettres soulignées par l’EDB et de contrôler les restitutions des textes décisifs (chapitre~\ref{ch:textes_decisifs}).
\item \emph{La collation des fastes.} La table consulaire doit être confrontée à la CLRE pour les années contestées.
\end{enumerate}

Ces conditions ne sont pas un ajout périphérique au livre. Elles déterminent les propositions que le lecteur peut tenir pour établies, celles qu’il peut tester à partir des données publiées, et celles qui restent à instruire.
